\begin{lstlisting}[
    language=Lisp,
    caption={Invariant and randomness mapping used in the theorem},
    label={lst:domino-theorem-equivalence-invariant},
    commentstyle=\small\rmfamily\itshape,
    basicstyle=\small\ttfamily,
    columns=fixed]
    (define-fun randomness-mapping-GARBLE
    (
        (sample-id-left SampleId)
        (sample-id-right SampleId)
        (sample-offset-left Int)
        (sample-offset-right Int)
    )
    Bool
    (or
        (and
            (= sample-id-left (sample-id "R" "GARBLE" "R_00"))
            (= sample-id-right (
                sample-id "SimFreeXOR0" 
                "GARBLE" "R_00")
            )
            (= sample-offset-left 0)
            (= sample-offset-right 0)
        )
        (and
            (= sample-id-left (
                sample-id "LINRKKDM" 
                "ENC" "delta")
            )
            (= sample-id-right (
                sample-id "SimFreeXOR0" 
                "GARBLE" "delta")
            )
            (= sample-offset-left 0)
            (= sample-offset-right 0)
        )
    )
    )
    (define-state-relation invariant
    (left-game right-game) 
    (and
        (= 
            (maybe-get left-game.GLINRKKDM.delta) 
            (right-game.SimFreeXOR0.delta))
    )
    )
\end{lstlisting}
